% ==============================================================================
% 249_FFGFT_Falsifizierbarkeit_En_ch.tex
% Chapter content – included via wrapper
% Author: Johann Pascher · May 2026
% ==============================================================================

\section*{Abstract}

This document is a companion to Dok.~248 (FFGFT --- Epistemological Position). It clarifies the philosophy-of-science position of FFGFT against metaphysical speculation: FFGFT is a physical theory with empirical content and concrete falsification conditions. The postulate of ``pure, unrestrained information'' as a primitive (Myser 2026) is, by contrast, a metaphysical framing without derivation power --- not falsifiable, not connectable to known physics. The decisive difference lies not only in the choice of primitive, but in the consequences of that choice.

\tableofcontents
\newpage

% ------------------------------------------------------------------------------
\section{Falsifiability as a Scientific Criterion}

A scientific theory distinguishes itself from metaphysical speculation through its \textbf{vulnerability to refutation} \cite{Popper1934}. A theory that cannot be falsified by any conceivable experiment makes no physical claims --- it makes claims about what we consider possible, not about what the world is.

This criterion is not a dogma but a methodological tool: it draws a clear line between theories that expose themselves to the world and those that evade it.

\begin{keyresult}[Popper's Criterion]
A theory is scientific if it is \textbf{principally falsifiable}: if there exist observations or experiments that would \emph{contradict} the theory and thereby refute it. A theory compatible with every possible observation says nothing about the world.
\end{keyresult}

% ------------------------------------------------------------------------------
\section{FFGFT Is Falsifiable}

FFGFT makes concrete, predictive claims that could in principle be refuted.

\subsection{Minimal Length Scale $L_0$}

\begin{equation}
  L_0 = \xipar \cdot \lP \approx 2.2 \times 10^{-39} \text{ m}
\end{equation}

This absolute lower bound on spacetime structure is a prediction of FFGFT. If future high-energy experiments or cosmological observations consistently showed structures \emph{below} this scale, FFGFT would be falsified. Conversely, if structures consistently remained at or above this bound, this supports the prediction.

\subsection{Fractal Dimension $D_f$}

The fractal dimension of spacetime in FFGFT is:
\begin{equation}
  D_f = 2.999867
\end{equation}
This is a deviation from exactly 3 by a factor of $1.33 \times 10^{-4}$. This deviation is in principle measurable via:
\begin{itemize}
  \item \textbf{Gamma-ray burst time-of-flight dispersion}: Energy-dependent light travel time from cosmological sources is an established search channel for quantum gravity effects. FFGFT predicts a specific, $\xipar$-dependent dispersion structure.
  \item \textbf{Gravitational wave phase structure}: Deviations in wave dispersion at extremely high frequencies.
\end{itemize}
If a measurement showed a deviation incompatible with $D_f = 2.999867$, FFGFT would be falsified.

\subsection{No Singularity in Black Holes}

FFGFT predicts that black holes contain no singularity --- because $L_0$ sets an absolute lower density bound that cannot be structurally violated.

This prediction is distinguishable from general relativity through:
\begin{itemize}
  \item \textbf{Gravitational wave echoes}: Quantum gravity models without singularities predict characteristic echo signals after mergers, potentially detectable by future generations of LIGO/Virgo.
  \item \textbf{Hawking radiation spectrum}: FFGFT predicts a modified spectrum reflecting the $T^4$ winding structure --- distinguishable from the purely thermal spectrum of standard theory.
\end{itemize}

\subsection{Coupling Constants from $\xipar$}

FFGFT derives all fundamental coupling constants from a single parameter $\xipar = 1.333 \times 10^{-4}$. Every experimental refinement of these constants is a test of FFGFT:
\begin{itemize}
  \item Fine structure constant $\alpha^{-1} = 137.036$ --- FFGFT derivation: $\alpha^{-1} \approx 137.036$ (deviation $< 0.001\%$)
  \item Anomalous magnetic moment of the electron $a_e$ --- FFGFT correction: $K_\text{frak}^{3/2}$ reduces deviation from $2.03\%$ to $0.014\%$
  \item Koide formula for lepton masses --- FFGFT provides a geometric derivation
\end{itemize}
Each of these derivations can be confirmed or refuted by more precise measurements.

% ------------------------------------------------------------------------------
\section{The Postulate of ``Pure Information'' Is Not Falsifiable}

Atticus Myser's postulate --- ``pure unrestrained information'' as primitive, without structure, without distinction, without binding --- makes \textbf{no predictive claims}.

\subsection{No Derivable Quantities}

From ``pure unrestrained information'' no measurable quantity follows, no length scale, no coupling constant, no field equation. The postulate describes what is suspected \emph{before} all structure --- but without structure there is no derivation, and without derivation there is no prediction.

\subsection{Compatible With Every Observation}

Since nothing follows from the postulate, it cannot be refuted by anything. Every observation --- whether it supports FFGFT, string theory, MOND, or standard cosmology --- is compatible with ``pure unrestrained information'', because the postulate excludes no observation.

\begin{wichtig}[Not Falsifiable = Not Physical]
A theory compatible with every observation makes no physical claims. The postulate of ``pure information'' is therefore not a physical theory --- it is a \textbf{metaphysical framing}. This is not inherently worthless, but it is a fundamentally different claim from that of FFGFT.
\end{wichtig}

\subsection{No Occlusion Mechanism Derivable}

Atticus's own theory of ``O-bits'' and sequential adhesion (occlusion) makes interesting structural claims. But these do not follow from the postulate of ``pure information'' --- they are introduced \emph{additionally}, as a separate rule. This rule is itself already structure. The postulate of the primitive does not carry the burden of occlusion --- it is imported from outside.

% ------------------------------------------------------------------------------
\section{The Philosophy-of-Science Distinction}

FFGFT and Atticus's approach differ not only in the choice of primitive ($T^4$ vs.\ ``pure information''), but in a deeper methodological point:

\begin{longtable}{>{\bfseries}p{3.5cm} p{5cm} p{5cm}}
\toprule
Criterion & FFGFT & ``Pure information'' (Myser) \\
\midrule
Primitive & $T^4 + \xipar$ (geometric, posited) & ``pure unrestrained information'' (posited) \\
Operational & Yes --- derivations, field equations, measurement & No --- no derivation possible \\
Falsifiable & Yes --- $L_0$, $D_f$, singularities, $\alpha$ & No --- no observation excluded \\
Predictions & Concrete, numerical & None \\
Connectable & To Standard Model, GR, QFT & Not directly \\
Type of claim & Physical theory & Metaphysical framing \\
\bottomrule
\end{longtable}

\begin{erkenntnis}[The Decisive Distinction]
\textbf{FFGFT is a physical theory with empirical content.} The postulate of ``pure information'' is a metaphysical speculation without derivation power. The decisive difference lies not only in the choice of axiom, but in the \emph{consequences} of that choice: what follows from the primitive? how much follows? and can what follows be refuted by observation?
\end{erkenntnis}

% ------------------------------------------------------------------------------
\section{Falsifiability and Epistemic Openness}

This distinction is not an attack on metaphysical framings. Metaphysical questions --- why is there something rather than nothing? what is the essence of information? what lies before all structure? --- are legitimate and important questions. They are simply not decidable by physical experiments.

FFGFT provides no answer to these questions. It leaves ``why $T^4$?'' explicitly open (cf.\ Dok.~248, Section 8). This is structural honesty: FFGFT describes \emph{from within}, with the means available to an inside description.

What distinguishes FFGFT from metaphysical speculation is not the absence of an answer to the ultimate question --- but the presence of concrete, refutable claims about the world we see.

\begin{erkenntnis}
\textbf{FFGFT is falsifiable but epistemically modest.} It makes no claims about what lies before $T^4$. It makes precise claims about what follows from $T^4 + \xipar$ --- and these claims can be refuted by observation. That is the claim of a physical theory, nothing more and nothing less.
\end{erkenntnis}

% ------------------------------------------------------------------------------
\section*{Summary}

\begin{longtable}{>{\bfseries}p{7.2cm} p{7.2cm}}
\toprule
Claim & Position \\
\midrule
FFGFT makes no predictive claims & False --- $L_0$, $D_f$, singularities, $\alpha$ are predictions \\
FFGFT is not falsifiable & False --- each prediction is in principle refutable \\
``Pure information'' is a physical theory & No --- metaphysical framing without derivation power \\
``Pure information'' is falsifiable & No --- no observation is excluded \\
Both approaches are scientifically equivalent & No --- only FFGFT is a physical theory in Popper's sense \\
Metaphysical framings are worthless & No --- they are legitimate, but of a different kind \\
FFGFT answers ``why $T^4$?'' & No --- this question remains explicitly open (cf.\ Dok.~248) \\
\bottomrule
\end{longtable}
